\documentclass{article}

\usepackage{amsmath,graphicx,epsfig}
% \usepackage{epstopdf}
% \usepackage{euscript}
\usepackage{amsfonts}
\usepackage{amssymb}
\usepackage{float}

\usepackage[framemethod=tikz]{mdframed}

\usepackage{setspace}
% \doublespacing
\usepackage[skip=10pt plus1pt, indent=20pt]{parskip}
% \parskip=1\baselineskip \advance\parskip by 0pt plus 2pt

% %%%%%%%%%%%%%%%%%%%%%%%%%%%%%%%%%%%%%%%%%%%%%%%%%%%%%%%%%%%%%%%%%%%%%%%%%%%%%%%%%%%%%%%%%%%%%%%%%%%%%%%%%%%%%%%%%
% We need hyperref for using LATEXDIFF
\usepackage{hyperref}

%DIF PREAMBLE EXTENSION ADDED BY LATEXDIFF
%DIF UNDERLINE PREAMBLE %DIF PREAMBLE
\RequirePackage[normalem]{ulem} %DIF PREAMBLE
\RequirePackage{color}\definecolor{RED}{rgb}{1,0,0}\definecolor{BLUE}{rgb}{0,0,1} %DIF PREAMBLE
\providecommand{\DIFaddtex}[1]{{\protect\color{blue}\uwave{#1}}} %DIF PREAMBLE
\providecommand{\DIFdeltex}[1]{{\protect\color{red}\sout{#1}}}                      %DIF PREAMBLE
%DIF SAFE PREAMBLE %DIF PREAMBLE
\providecommand{\DIFaddbegin}{} %DIF PREAMBLE
\providecommand{\DIFaddend}{} %DIF PREAMBLE
\providecommand{\DIFdelbegin}{} %DIF PREAMBLE
\providecommand{\DIFdelend}{} %DIF PREAMBLE
\providecommand{\DIFmodbegin}{} %DIF PREAMBLE
\providecommand{\DIFmodend}{} %DIF PREAMBLE
%DIF FLOATSAFE PREAMBLE %DIF PREAMBLE
\providecommand{\DIFaddFL}[1]{\DIFadd{#1}} %DIF PREAMBLE
\providecommand{\DIFdelFL}[1]{\DIFdel{#1}} %DIF PREAMBLE
\providecommand{\DIFaddbeginFL}{} %DIF PREAMBLE
\providecommand{\DIFaddendFL}{} %DIF PREAMBLE
\providecommand{\DIFdelbeginFL}{} %DIF PREAMBLE
\providecommand{\DIFdelendFL}{} %DIF PREAMBLE
%DIF HYPERREF PREAMBLE %DIF PREAMBLE
\providecommand{\DIFadd}[1]{\texorpdfstring{\DIFaddtex{#1}}{#1}} %DIF PREAMBLE
\providecommand{\DIFdel}[1]{\texorpdfstring{\DIFdeltex{#1}}{}} %DIF PREAMBLE
%DIF COLORLISTINGS PREAMBLE %DIF PREAMBLE
\RequirePackage{listings} %DIF PREAMBLE
\RequirePackage{color} %DIF PREAMBLE
\lstdefinelanguage{DIFcode}{ %DIF PREAMBLE
%DIF DIFCODE_UNDERLINE %DIF PREAMBLE
  moredelim=[il][\color{red}\sout]{\%DIF\ <\ }, %DIF PREAMBLE
  moredelim=[il][\color{blue}\uwave]{\%DIF\ >\ } %DIF PREAMBLE
} %DIF PREAMBLE
\lstdefinestyle{DIFverbatimstyle}{ %DIF PREAMBLE
	language=DIFcode, %DIF PREAMBLE
	basicstyle=\ttfamily, %DIF PREAMBLE
	columns=fullflexible, %DIF PREAMBLE
	keepspaces=true %DIF PREAMBLE
} %DIF PREAMBLE
\lstnewenvironment{DIFverbatim}{\lstset{style=DIFverbatimstyle}}{} %DIF PREAMBLE
\lstnewenvironment{DIFverbatim*}{\lstset{style=DIFverbatimstyle,showspaces=true}}{} %DIF PREAMBLE

\newcommand{\difDelContent}[1]{
  \DIFdelbegin \DIFdel{#1}\DIFdelend 
}
\newcommand{\difAddContent}[1]{
  \DIFaddbegin \DIFadd{#1}\DIFaddend 
}

%DIF END PREAMBLE EXTENSION ADDED BY LATEXDIFF
% %%%%%%%%%%%%%%%%%%%%%%%%%%%%%%%%%%%%%%%%%%%%%%%%%%%%%%%%%%%%%%%%%%%%%%%%%%%%%%%%%%%%%%%%%%%%%%%%%%%%%%%%%%%%%%%%%

% %%%%%%%%%%%%%%%%%%%%%%%%%%%%%%%%%%%%%%%%%%%%%%%%%%%%%%%%%%%%%%%%%%%%%%%%%%%%%%%%%%%%%%%%%%%%%%%%%%%%%%%%%%%%%%%%%
% commands for simplifying the boxes
\newcommand{\refereeOverallComment}[2]{
  \noindent \textbf{Referee #1 gives their overall opinion as follows:}
  \begin{mdframed}[linecolor=black]
  #2
  \end{mdframed}
}

\newcommand{\refereeNoOneBox}[2]{
  \begin{mdframed}[linecolor=black]
  \noindent \textbf{Referee 1: (#1)} #2
  \end{mdframed}
}

\newcommand{\refereeNoTwoBox}[2]{
  \begin{mdframed}[linecolor=black]
  \noindent \textbf{Referee 2: (#1)} #2
  \end{mdframed}
}

\newcommand{\mdframedBox}[1]{
  \begin{mdframed}[linecolor=blue,topline=false,bottomline=false]
  \noindent #1
  \end{mdframed}
}

\newcommand{\response}[2]{
  \noindent \textbf{Response to (#1):} #2
}
% %%%%%%%%%%%%%%%%%%%%%%%%%%%%%%%%%%%%%%%%%%%%%%%%%%%%%%%%%%%%%%%%%%%%%%%%%%%%%%%%%%%%%%%%%%%%%%%%%%%%%%%%%%%%%%%%%

% %%%%%%%%%%%%%%%%%%%%%%%%%%%%%%%%%%%%%%%%%%%%%%%%%%%%%%%%%%%%%%%%%%%%%%%%%%%%%%%%%%%%%%%%%%%%
% SI unit
\usepackage{siunitx}
\DeclareSIUnit\torr{torr}
\DeclareSIUnit\ppm{ppm}
\DeclareSIUnit\ppb{ppb}
\DeclareSIUnit\bar{bar}
\DeclareSIUnit\gauss{G}
\DeclareSIUnit\neV{\nano\electronvolt}
\DeclareSIUnit\peV{\pico\electronvolt}
\DeclareSIUnit\kms{km/s}
\sisetup{print-unity-mantissa=false}  % To display only the exponent (i.e., 10⁻¹¹) 
\sisetup{separate-uncertainty = false} % use parenthesis style for uncertainty
% \sisetup{uncertainty-mode = separate} % use +/- style for uncertainty
\sisetup{uncertainty-mode = compact-marker} % turn this on to use style like 123.45(1.20)
% otherwise we  turn this on to use style like 123.45(120)
% %%%%%%%%%%%%%%%%%%%%%%%%%%%%%%%%%%%%%%%%%%%%%%%%%%%%%%%%%%%%%%%%%%%%%%%%%%%%%%%%%%%%%%%%%%%%

\usepackage{changes}

% %%%%%%%%%%%%%%%%%%%%%%%%%%%%%%%%%%%%%%%%%%%%%%%%%%%%%%%%%%%%%%%%%%%%%%%%%%%%%%%%%%%%%%%%%%%%
% for chemical forumals
\usepackage{chemformula}
% %%%%%%%%%%%%%%%%%%%%%%%%%%%%%%%%%%%%%%%%%%%%%%%%%%%%%%%%%%%%%%%%%%%%%%%%%%%%%%%%%%%%%%%%%%%%


\begin{document}

\title{Response to Referee 2 for Manuscript ID universe-4542087}
%==================================

\date{\today}

\maketitle

\noindent Dear Editors of \textit{Universe} and the Referee,

% \noindent 
We would like to express our sincere gratitude to the Editors and the Referees. 
The referees understood the idea of this paper very well. They spotted a few important points and made useful comments and suggestions, which have helped improve our manuscript. 
Below we give the responses to the reviewers' comments and describe the related changes to the manuscript.



\refereeOverallComment{2}{This is an interesting and well organized manuscript.  The authors describe a
commissioning measurement in the CASPEr-Gradient-LF experiment to search for
couplings between the gradient of $\sim \SI{5.6}{\neV/c^2} $ axion field and proton spins in a
liquid methanol sample.  The measurements were run in a liquid helium cryostat
containing a superconducting solenoid with a bias field of 0.1 T.  In the
apparatus, an axion-induced signal would appear as a temporal precession in the
magnetization of the sample, measured with a readout chain that includes a
dc-SQUID and lock-in amplifier.  The sensitivity of the system is validated by
detection of RF test pulses that induce a magnetic field perpendicular to the
bias field.  The manuscript contains useful high-level information in the main
text, followed by detailed calculations of the readout, signal model, and
efficiency in the Appendices.
}

We are very glad to see the positive evaluation and constructive comments from the referee. Point-to-point responses are addressed as follows. 

\refereeNoOneBox{1}{The manuscript could probably benefit from more contextual information.  For
example, the authors highlight the applicability of the approach toward
``hitherto unexplored mass-coupling parameter space'', but without a direct
explanation.  In particular, the targeted mass range of 5.6 \unit{neV/c^2} is indeed
higher than in other searches such as in [50] (0.4 $->$ 4 \unit{feV/c^2}), or in the
NASDUCK experiment [e.g. Nature Communications volume 14, Article number: 5784
(2023)] (1.e-12 \unit{eV/c^2} $->$ 2e-10 \unit{eV/c^2}), but the reader is not informed directly
that the relatively high mass range here is accessible because of the
hyperpolarized sample being used.  If the latter is true, it would be helpful to
state this.  Table II hints at this information with predicted improvements
based on higher polarizations in liquid 129Xe, but otherwise the reader is left
to draw their own conclusions with respect to past measurements.
}

\response{1}{Thank you for pointing this out. 
We realize here more contextual information is necessary. 
There are two aspects in the mass-coupling parameter space: mass range and experimental sensitivity to the coupling.  
The mass range available for DM search using the CASPEr-Gradient apparatuses depends on the gyromagnetic ratio of the sample in use and the bias field. For example, with liquid xenon sample, atomic proton sample and a 0.1\,T magnet, we can explore a mass range of \SI{4.1}{\peV/c^2} to \SI{17}{\neV/c^2}.  
The sensitivity in that mass range depends on the polarization of the sample. With a hyperpolarized sample we plan to boost the sensitivity compared to thermally polarised sample by some orders of magnitude. 
We tried to indicate that, with hyperpolarized samples CASPEr will put the most stringent constraint in the accessible mass range.
Therefore, we modify the sentence in the abstract as: }
\mdframedBox{
\difDelContent{It opens the door to probing large swaths of hitherto unexplored mass-coupling parameter space in the future by using hyperpolarized samples.}

\difAddContent{
CASPEr-Gradient Low-Field will probe the mass range from 
$\SI{4.1}{peV/c^2}$ to $\SI{17}{neV/c^2}$
with hyperpolarized samples to boost the sensitivity beyond the astronomical limits.
}
% \DIFaddend 
}
We hope there is enough contextual information for the readers now.  

~\
~\

\refereeNoOneBox{2}{Similarly, the authors do not discuss the lower sensitivity achieved here with
respect to the other measurements as above ($\mathrm{g_{aP}}$~3.e-2/GeV vs. $\mathrm{g_{aNN}}$~1.e-10/GeV
[50] and $\mathrm{g_{aP}}$P~1.e-5/GeV [NASDUCK]).  Clearly the mass ranges are different, but
what are the technical limitations in this mass range compared with lower mass
ranges?  A reader might infer that the relaxation time T2 and T2*, quantified
empirically in this manuscript, might be the systematic effect that makes this
neV mass regime challenging.  Whereas the NASCUCK measurement operates
explicitly in the spin exchange relaxation free (SERF) regime.  Is that the key
technical challenge that makes hyperpolarized samples relevant for neV searches
such as CASPEr-Gradient-LF?  If not, then why are the hyperpolarized samples
needed here?
}

\response{2}{ 
It would be natural for readers to compare CASPEr to other experiments, and it is true that they deserve explanations on that.  
You are right about relaxation time T2 and T2* being systematic, frequency dependent effects that make this neV mass regime challenging.  
However, relaxation time is only one of several factors limiting our sensitivity. 
Some other factors are: 
\begin{itemize}
    \item polarization 
    \item axion coherence time 
    \item SQUID noise 
    \item spin projection noise (SPN) 
    \item Number of spins in sample 
\end{itemize}
The most significant technical limitation of our experiment that limited our sensitivity is the polarization of the sample. 
We want to use hyperpolarized samples to increase the sensitivity at a given mass range. Hyperpolarization increase the magnetization of the sample, thus improving the signal strength. This is true for experiments at any frequency/mass range.  
Currently we employ temperature control and delivery system to provide the CASPEr-gradient setup with liquid hyperpolarized xenon. 
To make this clear we added this text to the manuscript:  
\mdframedBox{As evident from Table I, the degree of polarization of the sample presents the largest limiting factor to CASPEr-Gradient-LF sensitivity. 
For that reason, in the future we will employ a hyperpolarized liquid Xe sample, which can boost sensitivity by up to several orders of magnitude alone.} 
} 
~\
~\



\refereeNoOneBox{3}{Other comments are:
The description of the data analysis is clear and useful.  The analysis consists
of filtering of the raw signals, followed by matched filtering according to the
expected axion signal lineshape.  I do not have recommendations for changes.
}

\response{3}{Thank you for your appreciation on the data analysis part. } 
~\
~\
\clearpage

\refereeNoOneBox{4}{Spin exchange optical pumping (SEOP) is referenced in Table II, but is not
defined until the conclusion.}

\response{4}{It was a mistake. Thank you again for the careful reading. Now we have included the information of the acronyms in the table as: }

\mdframedBox{
 Brute Force$^\text{c}$ PHIP$^\text{d}$ / DNP$^\text{e}$  SEOP$^\text{f}$
\begin{itemize}
    \footnotesize
    % \item[a] Here we take \SI{180}{\kelvin} methanol prepolarized at \SI{2}{\tesla} and measurement time $100\,\tau_a$ as the conditions for computing the baseline of future improvements. 
    % \item[b] \ch{^{129}Xe}-enriched sample. The \ch{^{129}Xe} concentration can be close to unity. The details of calculations about liquid \ch{^{129}Xe} sample can be found in Ref.\,\cite{kimball2018overview}.
    \item[c] Prepolarization in a strong magnetic field and/or at low temperature. 
    \item[d] Parahydrogen-Induced Polarization (PHIP). 
    \item[e] Dynamic Nuclear Polarization (DNP).
    \item[f] Spin Exchange Optical Pumping (SEOP).
\end{itemize}
}  
~\
~\

\refereeNoOneBox{5}{Table I:  What is spin projection noise, and why is it higher at 1\,kHz than at
4.3 MHz?
}

\response{5}{
We admit that the description of spin projection noise is missing from the text.
In section II we added a sentence to introduce spin-projection noise:
\mdframedBox{
\difAddContent{One intrinsic noise source is due to the spins within the
sample, known as spin-projection noise (SPN).}}
And in section III we added the following text to describe the dependence of spin-projection noise on frequency:
\mdframedBox{
\difAddContent{
The PSD of a full 100\,s data set is plotted in Fig.\,2\,(c).
The resulting noise floor of our PSD spectra is determined by spin-projection noise, SQUID noise, electronic noise (e.g. the amplifier noise) and noise created by mechanical vibrations. 
The SPN is proportional to the number of spins and inversely proportional to the NMR linewidth [$\Delta\nu=1/(\pi T_2^*)$]. Since the NMR linewidth is proportional to the frequency, SPN scales inversely with frequency (see Table\,I).}}
}
~\
~\
\refereeNoOneBox{6}{typo:  ``Appe'' $->$ Appendix.}

\response{6}{We appreciate your careful reading and helpful correction. It has been corrected. } 
~\
~\


% % %%%%%%%%%%%%%%%%%%%%%%%%%%%%%%%%%%%%%%%%%%%%%%%%%%%%%%%%%%%%%%%%%%%%%%%%%%%%%%%%%%%%%%%%%%%%
% \refereeNoOneBox{}{referee's comment here}

% \response{}{Thank you for pointing out the mistake in our manuscript. ??????? Thus we add the following text, indicated by blue. }

% \mdframedBox{content} 
% % %%%%%%%%%%%%%%%%%%%%%%%%%%%%%%%%%%%%%%%%%%%%%%%%%%%%%%%%%%%%%%%%%%%%%%%%%%%%%%%%%%%%%%%%%%%%


~\

~\


\refereeOverallComment{2}{This letter provides an overview of a search for axion-like particles
with the Cosmic Axion Spin Precession Experiment Gradient. The results
discussed in this paper were a proof-of-principle for the
experiment.The results detailed in this work set limits on the
axion-to-proton coupling for axion-like particle dark matter in a
narrow region about the 1.348570 MHz frequency range. This mass range
has not been well covered by lab experiments and CASPER-Gradient has
the potential to cover a large amount of new parameter space. I
recommend the manuscript be published on PRD. I have a few minor
clarification questions before the manuscript can be considered for
publication.
}

We sincerely appreciate your review and comments. Your comments have helped us improve the manuscript. 
% Point-to-point responses are addressed as follows. 


\refereeNoTwoBox{1}{What was the operating temperature of the experiment?}

\response{1}{This is a good question, as temperature plays an important role in our experiments. 
1. The magnets, coils, Nb shield and the SQUID were kept in liquid helium bath, which is at 4.2\,K so that they can be superconducting. 
2. The temperature of the sample was regulated by flowing nitrogen gas next to it. Throughout the measurements, it was kept constant around 180\,K, which is above the freezing point of methanol. This is also mentioned in the Appendix.}

% \mdframedBox{ccccccccccc\DIFdelbegin \DIFdel{as }\DIFdelend \DIFaddbegin \DIFadd{from }\DIFaddend} 


% %%%%%%%%%%%%%%%%%%%%%%%%%%%%%%%%%%%%%%%%%%%%%%%%%%%%%%%%%%%%%%%%%%%%%%%%%%%%%%%%%%%%%%%%%%%%
\refereeNoOneBox{2}{The letter details the use of pi pulses to determine the relaxation
time. Are there any uncertainties associated with these pulses?
}

\response{2}{We determined the parameters of the $\pi$-pulse by a pulse-duration sweep. Of course there is uncertainty in it, which is related to the step in the sweep. 
But when it comes to the T2 measurement, the uncertainty of pulse is relatively small compared to the uncertainty in spin-echo signal amplitudes due to the noise [as can be found in the Fig.2(c)]. Therefore, we do not consider the uncertainty of the pulse, but only the uncertainty of the spin-echo amplitude for the T2 value given in the manuscript. }

% \mdframedBox{content}  
~\
~\
% %%%%%%%%%%%%%%%%%%%%%%%%%%%%%%%%%%%%%%%%%%%%%%%%%%%%%%%%%%%%%%%%%%%%%%%%%%%%%%%%%%%%%%%%%%%%

% %%%%%%%%%%%%%%%%%%%%%%%%%%%%%%%%%%%%%%%%%%%%%%%%%%%%%%%%%%%%%%%%%%%%%%%%%%%%%%%%%%%%%%%%%%%%
\refereeNoOneBox{3}{The author details a series of veto procedures to account reduce the
number of outliers from X to X''. Following the veto procedure, were
the number of observed outliers within the statistically expected
number of false positives?
}

\response{3}{Yes, the numbers are consistent with the expected number of false positives. Indeed a statistical expectation value for the number of outliers after each veto can be calculated. For example in the case of the outliers on resonance, X', this would simply be the the expected outliers $\langle X \rangle$ multiplied by the fraction of the width of the resonant window to the width of the analysis window (approx. 140Hz/8kHz). We have added the expectation values into our table II, and hope it has become more clear that way. }
% \mdframedBox{[check Table II in the manuscript]} 
~\
~\

% %%%%%%%%%%%%%%%%%%%%%%%%%%%%%%%%%%%%%%%%%%%%%%%%%%%%%%%%%%%%%%%%%%%%%%%%%%%%%%%%%%%%%%%%%%%%

% %%%%%%%%%%%%%%%%%%%%%%%%%%%%%%%%%%%%%%%%%%%%%%%%%%%%%%%%%%%%%%%%%%%%%%%%%%%%%%%%%%%%%%%%%%%%
\refereeNoOneBox{4}{What was the primary factor limiting the bandwidth to 240 Hz. What are
the key upgrades that will expand that bandwidth?
}

\response{4}{The main limiting factor is the time it takes to scan the full range accessible by our experiment by ramping the bias field. In addition, our experiment has a limited duty cycle due to the need for manual intervention during each frequency scan, such as adjusting the magnetic field, initiating NMR diagnostics. In future iterations, this procedure will be fully automated, allowing for longer measurement times and thus broader bandwidth coverage.  }

% \mdframedBox{content}  
~\
~\
% %%%%%%%%%%%%%%%%%%%%%%%%%%%%%%%%%%%%%%%%%%%%%%%%%%%%%%%%%%%%%%%%%%%%%%%%%%%%%%%%%%%%%%%%%%%%

% %%%%%%%%%%%%%%%%%%%%%%%%%%%%%%%%%%%%%%%%%%%%%%%%%%%%%%%%%%%%%%%%%%%%%%%%%%%%%%%%%%%%%%%%%%%%
\refereeNoOneBox{5}{Fig 4- Would it be possible to include in your plot any limits set by
other experiments? Even though what is being presented is a
proof-of-concept, it would still be helpful to have additional context
on how the results compare?
}

\response{5}{
We have considered this and, in fact, have this plot available.
% \begin{figure}[h]
%     \centering
%     \includegraphics[width=\linewidth]{AxionProtoncouplinglimit.png}
%     \caption{The limit set by our commissioning measurement is depicted as ``CASPEr-gradient". Taken from \url{https://cajohare.github.io/AxionLimits/docs/app.html.}}
%     \label{fig:enter-label}
% \end{figure}

However, we would prefer not to include it in the current manuscript because it may give a false impression to the reader that CASPER-G-Low-Field does not possess competitive sensitivity. We could add our projections for the forthcoming run with a hyperpolarized sample, but this, again, may be misleading to readers expecting the actual results. Moreover, as you might be aware, Dr.\,Ciaran O'Hare operates a website where nearly all ultralight bosonic dark matter searches are combined in exclusion plots. This resource is used by many people in the field looking for the context that you ask for. We ensured that our latests results are included there, and feel that therefore it is not necessary to include it in our manuscript. Especially considering the proof-of-concept nature of the current work}

% \mdframedBox{content}  
% ~\
% ~\
% %%%%%%%%%%%%%%%%%%%%%%%%%%%%%%%%%%%%%%%%%%%%%%%%%%%%%%%%%%%%%%%%%%%%%%%%%%%%%%%%%%%%%%%%%%%%

We feel that the manuscript has significantly improved and hope you will find it acceptable for publication in the present form. 


% We hope that we have responded appropriately to all the sentences in the comments and made sufficient modifications. If there is any mistake or some comments are not well addressed, please let us know. We are happy to further improve the manuscript. 
\noindent Sincerely,



\noindent Julian Walter, Olympia Maliaka, Yuzhe Zhang and Dmitry Budker on behalf of the authors. 
%Pavel Fadeev and Derek F. Jackson Kimball on behalf of the authors.


\end{document} 